\scanpage{137}

% L1–6
Il faut rajouter le

% L2–6, left rule
\uline{Lemme} 1 Soit $F\subset E$, \deleted{loc. fact.}\inserted{$F\in\Delta_{\{d\}}(E,q)$}, Alors $F^\perp$ est loc. fact. direct et $\supset F$, et l'accouplement entre $F$ et $E/F^\perp$ est une dualité (donc $F=F^{\perp\perp}$, la formation de $F^\perp$ commute au chgt de base\ldots)

% L7–12
Il faut regarder \deleted{le noyau du \inserted{\uline{$F^\perp$}}} le composé $E\longrightarrow E^\vee\longrightarrow F^\vee$ et montrer que c'est \deleted{surjectif}\inserted{épi} -- le noyau étant $F^\perp$, le résultat s'ensuit. Or on est ramené au cas où $S$ est spectre d'un corps alg clos, [le seul cas non évident est celui où $\varphi$ \unsure{dégénérée|dégénère}, i.e. car $\unsure{k|h}=2$, $\operatorname{rg} E$ impair].

% L13–16
Mais \deleted{alors, si un $E\unsure{\text{---}|\to} E^\vee$ est le} si $a\in F$ est orth. à l'image de $E$, il est orth. à $E$ pour $\varphi$, \deleted{donc} ce qui joint à $q(\unsure{x|a})=0$ implique $a=0$, cqfd.

% L17–20
Reste à vérifier les assertions de connexion. Par dévissage, on est ramené, pour le cas \deleted{\unsure{$n$ p.}} $\operatorname{rg}E$ impair, ou $\operatorname{rg}E$ pair, $i_\nu<n$, \deleted{par}\inserted{le}

% lower left
% archive sheet number: 136

% L21–22
\uline{Lemme} \unsure{2|?} Soient $0\leq d<d'\leq n$, \inserted{\uline{NB} on admet ici $d=0$} \deleted{et supposons $d'<n$ si $\operatorname{rg}E$ pair}. Soit $F\in\Delta_{\{d\}}(E,q)$, consi%
